\documentclass[12pt,a4paper]{article}

% === Paquetes básicos ===
\usepackage[utf8]{inputenc}
\usepackage[spanish]{babel}
\usepackage{graphicx}
\usepackage{hyperref}
\usepackage{natbib}

% Márgenes
\usepackage[margin=2.5cm]{geometry}

% === Título ===
\title{Mi primer documento en \LaTeX}
\author{Tu Nombre}
\date{\today}

\begin{document}

\maketitle

\section{Introducción}

Este es un documento de ejemplo para practicar el flujo de trabajo en VS Code.
Lo único que necesitas recordar es:

\begin{itemize}
    \item \textbf{Ctrl+Alt+B} para compilar la primera vez.
    \item \textbf{Ctrl+Alt+V} para ver el PDF.
    \item \textbf{Ctrl+S} para guardar y recompilar automáticamente.
\end{itemize}

\section{Ejemplo de citas}

Según \citet{chacon2014pro}, Git es una herramienta fundamental para el control de versiones en proyectos de investigación. Otros autores también han explorado el uso de plataformas colaborativas en el ámbito académico \citep{perez2024github}.

\section{Conclusión}

Si puedes ver este PDF con las citas resueltas correctamente (no aparecen como [?]), tu instalación funciona perfectamente.

% === Bibliografía ===
\bibliographystyle{apalike}
\bibliography{biblio}

\end{document}
